% TOP_PROBLEM: {"id": 98, "title": "Sum-Free Subsets of [N]^d", "category_id": 2, "category": {"id": 2, "name": "combinatorics", "display_name": "Combinatorics", "description": "Counting problems, graph theory, discrete structures.", "slug": "combinatorics", "order_index": 2, "created_at": "2026-07-31T15:26:25.671Z"}, "status": "open"}
% FIRST_POSTED: null
% ENTRY_KIND: statement_only
% Imported from: unsolvedmath_additional_500.tex; UnsolvedMath ID 98
% !TeX program = lualatex
% Compile twice: lualatex 98.tex
% Self-contained: no external bibliography, images, or \input files.
% Numeric section labels and visible numbers are original UnsolvedMath IDs.
\documentclass[11pt,a4paper]{article}
\usepackage[margin=24mm,headheight=14pt]{geometry}
\usepackage{fontspec}
\setmainfont{Latin Modern Roman}
\setsansfont{Latin Modern Sans}
\setmonofont{Latin Modern Mono}
\usepackage{amsmath,amssymb,mathtools}
\usepackage{unicode-math}
\setmathfont{Latin Modern Math}
\usepackage{microtype}
\usepackage{enumitem,xurl,needspace}
\usepackage[dvipsnames]{xcolor}
\usepackage{hyperref}
\hypersetup{unicode=true,colorlinks=true,linkcolor=MidnightBlue,urlcolor=MidnightBlue,
 pdftitle={UnsolvedMath Problem 98},pdfauthor={Source-annotated compilation},bookmarksnumbered=true}
\usepackage{bookmark,tocloft,titlesec,fancyhdr}
\setlength{\cftsecnumwidth}{4.2em}
\cftsetpnumwidth{2.5em}
\cftsetrmarg{4em}
\titleformat{\section}{\Large\bfseries\raggedright}{\thesection}{0.7em}{}
\pagestyle{fancy}\fancyhf{}
\fancyhead[L]{\small\sffamily Additional UnsolvedMath statements}
\fancyhead[R]{\small Original catalogue IDs}
\fancyfoot[C]{\thepage}
\setlength{\parindent}{0pt}
\setlength{\parskip}{5pt plus 1pt}
\setlength{\emergencystretch}{3em}
\setcounter{tocdepth}{1}\setcounter{secnumdepth}{1}
\setlist[itemize]{leftmargin=3em,itemsep=4pt,topsep=4pt,parsep=0pt}
\allowdisplaybreaks
\newcommand{\N}{\mathbb{N}}
\newcommand{\Z}{\mathbb{Z}}
\newcommand{\Q}{\mathbb{Q}}
\newcommand{\R}{\mathbb{R}}
\newcommand{\C}{\mathbb{C}}
\newcommand{\F}{\mathbb{F}}
\newcommand{\A}{\mathbb{A}}
\newcommand{\PP}{\mathbb{P}}
\newcommand{\E}{\mathbb{E}}
\newcommand{\eps}{\varepsilon}
\newcommand{\dd}{\,\mathrm d}
\newcommand{\sref}[2]{\hyperlink{source:#1:#2}{[S#2]}}
\newif\ifshowcatalogue
\showcataloguetrue
\newcommand{\cataloguescope}[1]{\ifshowcatalogue
 \paragraph{Statement recorded in the supplied source.}
 #1\fi}
\begin{document}
% UnsolvedMath ID 98; source catalogue code GREEN-006

\setcounter{section}{97}

\section{Sum-Free Subsets of a Cartesian Grid}\label{98}

{\small\sffamily UnsolvedMath ID 98 · Combinatorics}

\subsection{Definitions and mathematical statement}

\paragraph{Shared notation.}Unless a section says otherwise, $\N=\{1,2,\ldots\}$,
$\Z$ is the ring of integers, $\Q,\R,\C$ are the rational, real, and complex
fields, $[N]=\{1,\ldots,\lfloor N\rfloor\}$, and $\log$ is the natural
logarithm. For real functions of a parameter tending to infinity,
$f\ll g$ and $f=O(g)$ mean $|f|\leq Cg$ eventually for some fixed $C$;
$f\gg g$ reverses this comparison; $f=o(g)$ means $f/g\to0$;
$f\sim g$ means $f/g\to1$; and $f\asymp g$ means both order bounds.
A subscript on $O$ or $\ll$ specifies allowed constant dependence.
The expression $n^{o(1)}$ in an upper bound means growth smaller than
$n^\epsilon$ for each fixed $\epsilon>0$. Local definitions govern any
exception, finite-field convention, or use of the same symbol for another
object. An asymptotic claim is not automatically an effective algorithm.



Write $[N]=\{1,\ldots,N\}$ and use coordinatewise addition in $\Z^d$. A set $A\subseteq[N]^d$ is sum-free if there are no $x,y,z\in A$ with $x+y=z$, including the case $x=y$. The extremal function is $s_d(N)=\max\{|A|:A\subseteq[N]^d,\ A\cap(A+A)=\varnothing\}$. The dimension $d$ is fixed while $N$ varies. \sref{98}{1} \sref{98}{2}

\paragraph{Statement recorded in the supplied source.}
Fix an integer $d$. What is the largest sum-free subset of $[N]^d$? \sref{98}{1}

\paragraph{Residual task reported by the archive.}

The embedded review (2026-08-17) records: Determine the maximum density for every fixed d. \sref{98}{1}

\subsection{Short English statement}

For a fixed dimension, how large can a subset of an integer grid be when adding two chosen points never produces another chosen point?

\subsection{Sources}

{\small

\begin{itemize}

\item[\hypertarget{source:98:1}{S1}] ULAM.AI, \emph{UnsolvedMath}, supplied \texttt{problems.json}, record 98 (GREEN-006), statement and background. The embedded literature-review date is 2026-08-17; that is the archive’s review date, not a fresh certification. Dataset landing page: \url{https://huggingface.co/datasets/ulamai/UnsolvedMath}.

\item[\hypertarget{source:98:2}{S2}] Ben Green, 100 Open Problems (author’s maintained problem list). \url{https://people.maths.ox.ac.uk/greenbj/papers/open-problems.pdf}. The author-hosted document was inspected on 27 September 2026; the archive’s local code is not a problem-number locator in this paper.

\item[\hypertarget{source:98:3}{S3}] S. Lepsveridze and Y. Sun, Size of the largest sum-free subset of [n]\textasciicircum{}3, [n]\textasciicircum{}4, and [n]\textasciicircum{}5, arXiv:2311.18289 (2023). \url{https://arxiv.org/abs/2311.18289}. Bibliographic information imported from the supplied record.

\item[\hypertarget{source:98:4}{S4}] S. Lepsveridze and Y. Sun, Size of the Largest Sum-Free Subset of [n]\textasciicircum{}3 and [n]\textasciicircum{}4, IMRN (2026). \url{https://doi.org/10.1093/imrn/rnag081}. Bibliographic information imported from the supplied record.

\end{itemize}

}

\label{end:98}
\end{document}
